\documentclass{article}
\usepackage[final]{neurips_2025}

\usepackage[utf8]{inputenc}
\usepackage[T1]{fontenc}
\usepackage{hyperref}
\usepackage{url}
\usepackage{microtype}
\usepackage{xcolor}
\usepackage{amsmath}

% --- Tighten spacing (NeurIPS-safe) ---
\usepackage{enumitem}
\setlist[itemize]{leftmargin=*,itemsep=0pt,topsep=0pt,parsep=0pt,partopsep=0pt}
\setlength{\parskip}{0pt}
\setlength{\abovedisplayskip}{2pt}
\setlength{\belowdisplayskip}{2pt}
\setlength{\textfloatsep}{4pt}
\setlength{\intextsep}{4pt}
\renewcommand{\baselinestretch}{0.97} % small squeeze; remove if disallowed

\title{Tri-Modal RAG VideoQA with Modality Routing}

\author{
  Aniket Salunke, Ketan More, Nazish Baliyan, Ritesh Thawkar \\
  Mohamed Bin Zayed University of Artificial Intelligence \\
  \texttt{\{aniket.salunke,ketan.more,nazish.baliyan,ritesh.thawkar\}@mbzuai.ac.ae}
}

\begin{document}
\maketitle
\vspace{-20pt}

\begin{abstract}\vspace{-2pt}
VideoQA evidence is distributed across subtitles, frames, and temporal segments, yet full-video processing remains prohibitively costly for real-world applications such as educational video search, media analysis, and accessible video assistants operating under strict latency and context constraints.
We propose a tri-modal retrieval-augmented pipeline with a modality router that dynamically allocates a fixed evidence budget across modalities per question, aiming to improve accuracy and reduce evidence mismatch compared to fixed retrieval policies.\vspace{-6pt}
\end{abstract}

\vspace{-4pt}
\paragraph{Problem statement.}
Given a video clip $V$ with aligned subtitles $S$ and a multiple-choice question $q$ with options $\{a_i\}_{i=1}^5$, predict the correct answer $\hat{a}$ using only a constrained evidence budget $B$.
A key challenge is \emph{modality dependence}: dialogue-heavy questions need subtitle evidence, appearance/object questions need frames, and action questions need temporal segments.
\textbf{Hypothesis:} routing the evidence budget across modalities based on the question improves end-to-end VideoQA compared to single-modality or fixed-budget retrieval.

\vspace{-2pt}
\paragraph{Related work.}
TVQA~\cite{lei2018tvqa} introduced multimodal VideoQA with aligned subtitles, establishing the benchmark we build on, while CLIP~\cite{radford2021clip} and VideoCLIP~\cite{xu2021videoclip} provide strong cross-modal encoders for frame- and segment-level retrieval.
Retrieval-augmented generation (RAG)~\cite{lewis2020rag} has been extended to video: Video-RAG~\cite{luo2024videorag} retrieves visually-aligned evidence for long videos, VRAG~\cite{gia2025vrag} targets long-form VideoQA, Sagare et al.~\cite{sagare2024videorag} scale context size, and Jeong et al.~\cite{jeong2025videorag} retrieve over full video corpora; yet UniversalRAG~\cite{yeo2025universalrag} highlights persistent modality gaps in unified-space retrieval.
None of these works dynamically \emph{route} a fixed evidence budget across modalities conditioned on the question type---the gap our project addresses.

\vspace{-2pt}
\paragraph{Dataset \& evaluation.}
\textbf{TVQA} (primary)~\cite{lei2018tvqa}: video + subtitles + (question, 5 answers).
We output $(\hat{a},E)$ where $E$ contains cited subtitle spans and timestamps plus selected keyframes/segments.
\textbf{Metrics:} QA accuracy; retrieval quality (Recall@$K$/MRR over evidence items, and temporal overlap w.r.t.\ annotated moments when available); evidence efficiency (items/tokens per query); and calibration/robustness checks (e.g., accuracy vs.\ budget $B$ and perturbed/noisy subtitles).

\vspace{-2pt}
\paragraph{Methodology.}
We build three evidence stores: (1) subtitle chunks (timestamped), (2) keyframes, and (3) short segments (2--4s windows).
We embed each store with pretrained encoders (text; image-language such as CLIP~\cite{radford2021clip}; video-text representations~\cite{xu2021videoclip,luo2021clip4clip}) and retrieve top-$K_m$ items per modality using dense retrieval~\cite{karpukhin2020dpr} with ANN indexing~\cite{johnson2017faiss}.
Our contribution is a \textbf{modality router} that predicts $(K_{\text{text}},K_{\text{img}},K_{\text{vid}})$ under total budget $B$, and retrieves accordingly from modality-specific indexes (avoiding unified-space modality gaps in multimodal RAG~\cite{yeo2025universalrag}).
Router training will follow: (A) supervision from question patterns/keywords (dialogue vs visual vs action), and/or (B) imitation of the best split found on a validation set.
Answer scoring uses evidence-conditioned features for each option (e.g., max/mean similarity to retrieved evidence per modality) and an MLP classifier to select $\hat{a}$ while returning top evidence as citations.

\vspace{-2pt}
\paragraph{Baselines.}
(i) \textbf{No-RAG}: answer using fixed context (uniform keyframes/segments + truncated subtitles), without any learned retrieval;
(ii) \textbf{Single-modality RAG}: retrieve from only one modality (text-only / frame-only / segment-only) under the same total budget $B$ to isolate modality contributions;
(iii) \textbf{Fixed tri-modal RAG}: retrieve from all modalities with a constant split of $B$ (uniform $(K_{\text{text}},K_{\text{img}},K_{\text{vid}})$), then score options using the same answer module as our method.

\vspace{-2pt}
\paragraph{Ablations.}
(a) \textbf{Routing}: router off vs on; keyword-supervised vs validation-imitation training; sensitivity to routing errors (forced wrong split);
(b) \textbf{Budget \& granularity}: sweep $B$; segment window (2s/4s) and frame rate; subtitle chunk size; compare uniform vs scene-change keyframes;
(c) \textbf{Ranking/scoring}: rerank on/off; similarity-only scoring vs MLP option scorer; encoder choice per modality and late-fusion strategies.
\textbf{Analysis:} per-question-type breakdown, evidence-mismatch audit (retrieval failure vs reasoning/scoring failure), and accuracy--compute trade-offs (latency vs accuracy and evidence budget vs accuracy).


\vspace{-2pt}
\paragraph{Work division.}
Work will be split by components with shared integration: a designated member leads preprocessing; another leads embedding/indexing + retrieval baselines; another leads router training + ablations; another leads option scoring + evaluation/demo. All members jointly define protocols, review results, and write the report/presentation.

\vspace{-2pt}
\paragraph{Software \& compute.}
Deliverables: reproducible preprocessing, indexing/retrieval, router, scorer, evaluation, and a demo script (question $\rightarrow$ answer + cited evidence) in a private GitHub repo.
Stack: PyTorch (ROCm), HuggingFace, OpenCLIP/CLIP, FAISS, decord/PyAV, NumPy/scikit-learn.
Compute: 1 GPU baseline; up to 8 AMD GPUs for parallel feature extraction and sweeps.

\vspace{-2pt}
\paragraph{Papers to read.}
\begin{itemize}
  \item \textbf{TVQA: Localized, Compositional Video Question Answering}~\cite{lei2018tvqa}: defines the dataset, task format, and evaluation protocol adopted in this project.
  \item \textbf{Video-RAG: Visually-Aligned Retrieval-Augmented Long Video Comprehension}~\cite{luo2024videorag}: the closest prior work on retrieval-augmented video understanding, which informs our indexing and retrieval design.
  \item \textbf{Learning Transferable Visual Models From Natural Language Supervision (CLIP)}~\cite{radford2021clip}: provides the backbone encoder for frame-level dense features central to our embedding pipeline.
  \item \textbf{UniversalRAG: Retrieval-Augmented Generation Over Multiple Corpora With Diverse Modalities and Granularities}~\cite{yeo2025universalrag}: exposes modality gaps in unified-space multimodal RAG, motivating our design of separate per-modality indexes.
  \item \textbf{Smart Routing for Multimodal Video Retrieval: When to Search What (ModaRoute)}~\cite{delarosa2025modaroute}: introduces LLM-based modality routing for video retrieval, directly related to our question-conditioned routing mechanism.
\end{itemize}

% References can go beyond one page.
\begin{thebibliography}{99}\setlength{\itemsep}{0pt}

\bibitem{lei2018tvqa}
J.~Lei, L.~Yu, T.~Bansal, and T.~Berg.
\newblock TVQA: Localized, compositional video question answering.
\newblock In \emph{Proceedings of EMNLP}, 2018.

\bibitem{radford2021clip}
A.~Radford, J.~W. Kim, C.~Hallacy, A.~Ramesh, G.~Goh, S.~Agarwal, G.~Sastry,
A.~Askell, P.~Mishkin, J.~Clark, G.~Krueger, and I.~Sutskever.
\newblock Learning transferable visual models from natural language supervision.
\newblock In \emph{Proceedings of ICML}, 2021.

\bibitem{xu2021videoclip}
H.~Xu, G.~Ghosh, P.-Y. Huang, D.~Okhonko, A.~Aghajanyan, F.~Metze,
L.~Zettlemoyer, and C.~Feichtenhofer.
\newblock VideoCLIP: Contrastive pre-training for zero-shot video-text understanding.
\newblock In \emph{Proceedings of EMNLP}, 2021.

\bibitem{luo2021clip4clip}
H.~Luo, L.~Ji, B.~Shi, H.~Huang, N.~Duan, T.~Li, X.~Chen, and M.~Zhou.
\newblock CLIP4Clip: An empirical study of CLIP for end-to-end video clip retrieval.
\newblock arXiv:2104.08860, 2021.

\bibitem{lewis2020rag}
P.~Lewis, E.~Perez, A.~Piktus, F.~Petroni, V.~Karpukhin, N.~Goyal, H.~K{\"u}ttler,
M.~Lewis, W.-t. Yih, T.~Rockt{\"a}schel, S.~Riedel, and D.~Kiela.
\newblock Retrieval-augmented generation for knowledge-intensive NLP tasks.
\newblock In \emph{Proceedings of NeurIPS}, 2020.

\bibitem{karpukhin2020dpr}
V.~Karpukhin, B.~O{\u{g}}uz, S.~Min, P.~Lewis, L.~Wu, S.~Edunov, D.~Chen, and W.-t. Yih.
\newblock Dense passage retrieval for open-domain question answering.
\newblock In \emph{Proceedings of EMNLP}, 2020.

\bibitem{johnson2017faiss}
J.~Johnson, M.~Douze, and H.~J{\'e}gou.
\newblock Billion-scale similarity search with GPUs.
\newblock arXiv:1702.08734, 2017.

\bibitem{luo2024videorag}
Y.~Luo, X.~Zheng, G.~Li, S.~Yin, H.~Lin, C.~Fu, J.~Huang, J.~Ji, F.~Chao, J.~Luo, and R.~Ji.
\newblock Video-RAG: Visually-aligned retrieval-augmented long video comprehension.
\newblock arXiv:2411.13093, 2024. (Accepted at NeurIPS 2025)

\bibitem{gia2025vrag}
B.~T. Gia, K.~Le, T.~Do, T.-D. Mai, T.~D. Ngo, D.-D. Le, and S.~Satoh.
\newblock VRAG: Retrieval-augmented video question answering for long-form videos.
\newblock In \emph{CVPR Workshops (IViSE)}, 2025.

\bibitem{sagare2024videorag}
S.~R. Sagare, P.~Ullegaddi, N.~K.~S., N.~R., K.~Sarabhai, and R.~K.~S.~A.
\newblock VideoRAG: Scaling the context size and relevance for video question-answering.
\newblock In \emph{Proceedings of INLG (System Demonstrations)}, 2024.

\bibitem{jeong2025videorag}
S.~Jeong, K.~Kim, J.~Baek, and S.~J. Hwang.
\newblock VideoRAG: Retrieval-augmented generation over video corpus.
\newblock In \emph{Findings of ACL}, 2025.

\bibitem{yeo2025universalrag}
W.~Yeo, K.~Kim, S.~Jeong, J.~Baek, and S.~J. Hwang.
\newblock UniversalRAG: Retrieval-augmented generation over multiple corpora with diverse modalities and granularities.
\newblock arXiv:2504.20734, 2025.

\bibitem{delarosa2025modaroute}
K.~Dela~Rosa.
\newblock Smart routing for multimodal video retrieval: When to search what.
\newblock In \emph{ICCV Workshops (Multimodal Representation and Retrieval)}, 2025.

\end{thebibliography}

\end{document}
